% ============================================================
%  Chapter 1, Introduction
% ============================================================
\chapter{Introduction}
\label{ch:intro}

Time series forecasting is one of the most studied areas of applied machine
learning, partly because real-world data is almost always noisy.
Sensor dropouts, aggregation artefacts, and market microstructure noise all
degrade the signal-to-noise ratio, and models frequently struggle with this.
In the long-term forecasting (LTSF) setting the problem is harder still,
since a lookback window of length $L$ must support a prediction of the next $H$ steps,
where $H$ can exceed $L$~\cite{zhou2021informer}.

\section{Motivation}

Transformer-based models~\cite{zhou2021informer,wu2021autoformer} and the
Mamba SSM~\cite{gu2021s4,gu2023mamba} address data-quality issues through
implicit generalisation, relying on the model to learn around the noise.
This works on clean data, but there is no guarantee that the robustness
transfers to noisy regimes.

The Kalman filter~\cite{kalman1960new} takes a different route: it builds an
\emph{explicit} model of process noise ($Q$) and measurement noise ($R$), and
derives a gain $K$ from their ratio to decide how much to trust the current
measurement versus the prior state.
The \textbf{KCMamba} architecture presented in this thesis is built on
this principle. That inductive bias, if embedded in a neural network,
may help in noisy conditions.

\section{Research Questions}

Three questions are investigated: (1)~is KCMamba competitive on clean data
against standard LTSF models; (2)~how much MSE degradation do models suffer
under synthetically injected noise ($\sigma \in \{0, 1, 3, 5\}$);
and (3)~which architecture generalises better under scarce training data
(stride $\in \{1, 16\}$).

\section{Contributions}

The thesis makes five concrete contributions.
The \textbf{KCMamba} architecture embeds learned $Q_\text{net}$,
$R_\text{net}$, $K_\text{net}$ subnetworks inside a differentiable Kalman
filter, with a parallel prefix-scan GPU implementation for the recursion.
Measurements show that the learned $K$ genuinely adapts to noise:
$K \approx 0{.}71$ at $\sigma=0$ (passive pass-through) and
$K \approx 0{.}07$ at $\sigma=5$ (strong filtering).

Benchmark results on three datasets (Exchange-Rate, ETTh1, ETTh2) show
that Fair Mamba degrades by $28\,234\%$ under noise, whereas KCMamba
degrades by only $1\,528\%$ ($18.5\times$ smaller degradation).
The long-sequence analysis ($T \in \{96,192,336,512\}$, ETTm1 and
Exchange) confirms KCMamba's advantage at longer contexts:
at $T{=}336$, $\sigma{=}1$ it achieves $31.7\%$ lower MSE than LSTM
($0.502$ vs $0.735$); at $T{=}512$, $\sigma{=}3$, KCMamba scores $0.773$
while Fair Mamba scores $2.121$. On ETTm1 at $\sigma{=}3$ KCMamba is
the only model whose MSE monotonically decreases with sequence length
($0.673 \to 0.466$, $-31\%$, $T{:}\,96 \to 512$).
In the stride $\times$ sequence-length sweep ($s \in \{1,16\}$,
$T \in \{192,336,512\}$), at $s{=}16$, $\sigma{\geq}3$,
Fair Mamba's MSE rises to $8.1$--$8.5$ while KCMamba stays at $1.1$--$1.3$.
All experiments are accompanied by an open-source Python/PyTorch/CUDA
implementation integrated into an asynchronous live-trading system.

\section{Thesis Structure}

Chapter~\ref{ch:irodalom} reviews the Kalman filter, SSMs, and the
LTSF literature.
Chapter~\ref{ch:architektura} describes the KCMamba architecture and
its mathematical underpinnings.
Chapter~\ref{ch:kiserletek} presents the experimental setup and results.
Chapter~\ref{ch:implementacio} covers the implemented software system.
Chapter~\ref{ch:osszegzes} draws conclusions and discusses future work.
